
\documentclass[11pt,letterpaper]{article}

% =========================================
% PAQUETES
% =========================================
\usepackage[utf8]{inputenc}
\usepackage[T1]{fontenc}
\usepackage[spanish]{babel}
\usepackage{lmodern}
\usepackage{amsmath,amssymb}
\usepackage{geometry}
\usepackage{xcolor}
\usepackage{tikz}
\usepackage{tcolorbox}
\usepackage{fancyhdr}
\usepackage{enumitem}
\usepackage{array}
\usepackage{multicol}
\usepackage{lastpage}

% TikZ antes de begin document
\usetikzlibrary{
 arrows.meta,
 calc,
 angles,
 quotes,
 positioning
}

\tcbuselibrary{breakable,skins}

% =========================================
% CONFIGURACIÓN
% =========================================
\geometry{
 left=1.8cm,
 right=1.8cm,
 top=2.2cm,
 bottom=2cm
}

\setlength{\parindent}{0pt}
\setlength{\parskip}{5pt}
\setlength{\headheight}{15pt}

% =========================================
% COLORES
% =========================================
\definecolor{azul}{RGB}{20,65,115}
\definecolor{celeste}{RGB}{232,243,252}
\definecolor{verde}{RGB}{0,130,100}
\definecolor{morado}{RGB}{110,50,160}
\definecolor{gris}{RGB}{246,248,251}
\definecolor{naranja}{RGB}{230,110,35}

% =========================================
% ENCABEZADO
% =========================================
\pagestyle{fancy}
\fancyhf{}

\lhead{\textcolor{azul}{
\textbf{Física -- Educación Diversificada}
}}

\rhead{\textcolor{azul}{J. Arias}}

\cfoot{
Página \thepage\ de \pageref{LastPage}
}

\renewcommand{\headrulewidth}{0.5pt}

% =========================================
% ESTILOS DE RECUADROS
% =========================================
\newtcolorbox{seccion}[1]{
 enhanced,
 breakable,
 colback=celeste,
 colframe=azul,
 colbacktitle=azul,
 coltitle=white,
 title=\textbf{#1},
 fonttitle=\bfseries,
 arc=2mm,
 boxrule=0.8pt
}

\newtcolorbox{ejercicio}[1]{
 enhanced,
 breakable,
 colback=white,
 colframe=azul!55,
 title=\textbf{Ejercicio #1},
 colbacktitle=celeste,
 coltitle=azul,
 fonttitle=\bfseries,
 arc=2mm,
 boxrule=0.6pt
}

\newtcolorbox{reto}{
 enhanced,
 breakable,
 colback=orange!5,
 colframe=naranja,
 title=\textbf{Desafío},
 colbacktitle=naranja,
 coltitle=white,
 arc=2mm
}

% =========================================
% COMANDOS
% =========================================
\newcommand{\respuesta}{
\vspace{0.3cm}
\textcolor{morado}{\textbf{Respuesta:}}
\hrulefill
\vspace{0.3cm}
}

\newcommand{\espacio}[1]{
\vspace{#1}
}

\newcommand{\casilla}{
$\square$
}

% Plano cartesiano reutilizable
\newcommand{\plano}{
\begin{tikzpicture}[scale=0.55]
\draw[step=1,gray!25,very thin]
(-6,-5) grid (6,5);

\draw[-{Stealth},thick]
(-6.3,0)--(6.4,0)
node[right] {$x$};

\draw[-{Stealth},thick]
(0,-5.3)--(0,5.4)
node[above] {$y$};

\foreach \x in {-6,-4,-2,2,4,6}{
\node[below,font=\scriptsize]
at (\x,0) {\x};
}

\foreach \y in {-4,-2,2,4}{
\node[left,font=\scriptsize]
at (0,\y) {\y};
}
\end{tikzpicture}
}

% =========================================
% DOCUMENTO
% =========================================
\begin{document}

% =========================================
% PORTADA
% =========================================
\begin{center}

{\Large\bfseries
\textcolor{azul}{
PRÁCTICA DE FÍSICA
}}

\vspace{0.15cm}

{\large
Magnitudes escalares y vectoriales
}

\vspace{0.2cm}

\textcolor{verde}{
\textbf{Operaciones y componentes rectangulares}
}

\end{center}

\begin{tcolorbox}[
 colback=gris,
 colframe=azul,
 arc=2mm
]

\textbf{Estudiante:}
\hrulefill

\vspace{0.3cm}

\textbf{Sección:}
\underline{\hspace{3cm}}
\hfill
\textbf{Fecha:}
\underline{\hspace{3cm}}

\end{tcolorbox}

\begin{tcolorbox}[
 colback=celeste,
 colframe=azul,
 title=Indicaciones
]

\begin{itemize}[leftmargin=*]
\item Resuelva cada ejercicio mostrando
todos los procedimientos.

\item Utilice calculadora científica
en modo grados.

\item Exprese las respuestas con
las unidades correspondientes.

\item Redondee a dos decimales
cuando sea necesario.

\item Realice las representaciones
gráficas solicitadas.
\end{itemize}

\end{tcolorbox}

% =========================================
% SECCIÓN I
% =========================================
\section*{I. Magnitudes escalares y vectoriales}

\begin{seccion}{
Clasificación de magnitudes
}

Marque con una X la clasificación
correcta de cada magnitud.

\begin{center}
\renewcommand{\arraystretch}{1.6}

\begin{tabular}{|p{8cm}|c|c|}
\hline
\textbf{Magnitud}
& \textbf{Escalar}
& \textbf{Vectorial}\\
\hline
Masa de $12\,\mathrm{kg}$
& $\square$ & $\square$\\
\hline
Fuerza de $50\,\mathrm{N}$ al norte
& $\square$ & $\square$\\
\hline
Tiempo de $35\,\mathrm{s}$
& $\square$ & $\square$\\
\hline
Velocidad de $20\,\mathrm{m/s}$ al este
& $\square$ & $\square$\\
\hline
Distancia de $250\,\mathrm{m}$
& $\square$ & $\square$\\
\hline
Aceleración de $4\,\mathrm{m/s^2}$
hacia abajo
& $\square$ & $\square$\\
\hline
Temperatura de $300\,\mathrm{K}$
& $\square$ & $\square$\\
\hline
Desplazamiento de $80\,\mathrm{m}$
al oeste
& $\square$ & $\square$\\
\hline
Rapidez de $15\,\mathrm{m/s}$
& $\square$ & $\square$\\
\hline
Fuerza de $120\,\mathrm{N}$
a $35^\circ$
& $\square$ & $\square$\\
\hline
\end{tabular}
\end{center}

\end{seccion}

% =========================================
% SECCIÓN II
% =========================================
\newpage

\section*{II. Representación gráfica}

\begin{seccion}{
Representación de vectores
}

Utilice los planos cartesianos
para representar los vectores.

\end{seccion}

\begin{ejercicio}{1}

Represente los siguientes vectores
utilizando una escala adecuada.

\begin{enumerate}[label=\alph*)]
\item $\vec A=40\,\mathrm{N}$ hacia el este.
\item $\vec B=30\,\mathrm{N}$ hacia el norte.
\item $\vec C=50\,\mathrm{N}$ hacia el oeste.
\item $\vec D=20\,\mathrm{N}$ hacia el sur.
\end{enumerate}

\begin{center}
\plano
\end{center}

\end{ejercicio}

\begin{ejercicio}{2}

Represente gráficamente los vectores:

\begin{enumerate}[label=\alph*)]
\item $\vec A=50\,\mathrm{N}$ a $30^\circ$.
\item $\vec B=40\,\mathrm{N}$ a $120^\circ$.
\item $\vec C=60\,\mathrm{N}$ a $210^\circ$.
\item $\vec D=35\,\mathrm{N}$ a $300^\circ$.
\end{enumerate}

Los ángulos se miden desde
el eje $x$ positivo.

\begin{center}
\plano
\end{center}

\end{ejercicio}

% =========================================
% SECCIÓN III
% =========================================
\newpage

\section*{III. Suma y resta gráfica}

\begin{seccion}{
Método del triángulo y paralelogramo
}

Resuelva mediante el método gráfico
indicado en cada ejercicio.

\end{seccion}

\begin{ejercicio}{3}

Dos desplazamientos tienen las
siguientes características:

\[
A=40\,\mathrm{m}
\quad\text{hacia el este}
\]

\[
B=30\,\mathrm{m}
\quad\text{hacia el norte}
\]

\begin{enumerate}[label=\alph*)]
\item Dibuje los vectores consecutivos.
\item Trace la resultante.
\item Determine gráficamente
su magnitud y dirección.
\end{enumerate}

Escala sugerida:
$1\,\mathrm{cm}=10\,\mathrm{m}$.

\begin{center}
\plano
\end{center}

\respuesta

\end{ejercicio}
\newpage
\begin{ejercicio}{4}

Dos fuerzas concurrentes actúan
sobre un cuerpo:

\[
F_1=60\,\mathrm{N}
\quad\text{a }0^\circ
\]

\[
F_2=40\,\mathrm{N}
\quad\text{a }60^\circ
\]

\begin{enumerate}[label=\alph*)]
\item Dibuje ambas fuerzas.
\item Complete el paralelogramo.
\item Trace la resultante.
\item Estime su magnitud y dirección.
\end{enumerate}

\begin{center}
\plano
\end{center}

\respuesta

\end{ejercicio}

\newpage

\begin{ejercicio}{5}

Considere los vectores:

\[
\vec A=50\,\mathrm{N}
\quad\text{hacia el este}
\]

\[
\vec B=30\,\mathrm{N}
\quad\text{hacia el norte}
\]

Represente gráficamente:

\[
\vec R=\vec A-\vec B
\]

Determine la magnitud y dirección
de la resultante.

\begin{center}
\plano
\end{center}

\respuesta

\end{ejercicio}

% =========================================
% SECCIÓN IV
% =========================================
\newpage

\section*{IV. Vectores perpendiculares}

\begin{seccion}{
Magnitud y dirección de la resultante
}

Determine la magnitud y dirección
de la resultante de cada par
de vectores perpendiculares.

\end{seccion}

\begin{ejercicio}{6}

Un estudiante camina
$12\,\mathrm{m}$ hacia el este
y después $16\,\mathrm{m}$
hacia el norte.

Determine:

\begin{enumerate}[label=\alph*)]
\item Magnitud del desplazamiento.
\item Dirección respecto al este.
\end{enumerate}

\espacio{4cm}

\respuesta

\end{ejercicio}

\begin{ejercicio}{7}

Sobre un objeto actúan dos fuerzas:

\[
F_1=45\,\mathrm{N}
\quad\text{hacia el este}
\]

\[
F_2=60\,\mathrm{N}
\quad\text{hacia el norte}
\]

Determine la magnitud y dirección
de la fuerza resultante.

\espacio{4cm}

\respuesta

\end{ejercicio}

\begin{ejercicio}{8}

Un vehículo se desplaza
$80\,\mathrm{m}$ hacia el oeste
y luego $60\,\mathrm{m}$
hacia el sur.

Determine:

\begin{enumerate}[label=\alph*)]
\item Magnitud del desplazamiento.
\item Dirección de la resultante.
\item Cuadrante correspondiente.
\end{enumerate}

\espacio{4cm}

\respuesta

\end{ejercicio}

% =========================================
% SECCIÓN V
% =========================================
\newpage

\section*{V. Componentes rectangulares}

\begin{seccion}{
Descomposición de vectores
}

Calcule las componentes horizontal
y vertical de cada vector.

\end{seccion}

\begin{ejercicio}{9}

Una fuerza de $100\,\mathrm{N}$
forma un ángulo de $35^\circ$
con el eje $x$ positivo.

Determine:

\[
F_x=\underline{\hspace{4cm}}
\]

\[
F_y=\underline{\hspace{4cm}}
\]

\espacio{3cm}

\end{ejercicio}

\begin{ejercicio}{10}

Un vector de magnitud
$150\,\mathrm{N}$ forma
un ángulo de $125^\circ$
con el eje $x$ positivo.

Determine sus componentes.

\espacio{4cm}

\respuesta

\end{ejercicio}
\newpage
\begin{ejercicio}{11}

Una fuerza de $200\,\mathrm{N}$
forma un ángulo de $230^\circ$
con el eje $x$ positivo.

Determine:

\begin{enumerate}[label=\alph*)]
\item Componente horizontal.
\item Componente vertical.
\item Cuadrante correspondiente.
\end{enumerate}

\espacio{4cm}

\respuesta

\end{ejercicio}

\begin{ejercicio}{12}

Un vector de $80\,\mathrm{N}$
forma un ángulo de $315^\circ$.

Determine sus componentes
rectangulares.

\espacio{4cm}

\respuesta

\end{ejercicio}

% =========================================
% SECCIÓN VI
% =========================================
\newpage

\section*{VI. Fuerzas concurrentes}

\begin{seccion}{
Resultantes por componentes rectangulares
}

Complete las tablas de componentes
y determine la magnitud y dirección
de cada resultante.

\end{seccion}

\begin{ejercicio}{13}

Sobre un cuerpo actúan tres fuerzas:

\[
F_1=100\,\mathrm{N}
\quad\text{hacia el este}
\]

\[
F_2=80\,\mathrm{N}
\quad\text{hacia el norte}
\]

\[
F_3=40\,\mathrm{N}
\quad\text{hacia el oeste}
\]

Complete:

\begin{center}
\renewcommand{\arraystretch}{1.8}

\begin{tabular}{|c|c|c|}
\hline
Fuerza & $F_x$ (N) & $F_y$ (N)\\
\hline
$F_1$ & &\\
\hline
$F_2$ & &\\
\hline
$F_3$ & &\\
\hline
$\sum$ & &\\
\hline
\end{tabular}
\end{center}

\espacio{3cm}

\respuesta

\end{ejercicio}

\newpage

\begin{ejercicio}{14}

Cuatro fuerzas actúan sobre
un mismo punto:

\[
F_1=120\,\mathrm{N},\quad 0^\circ
\]

\[
F_2=90\,\mathrm{N},\quad 60^\circ
\]

\[
F_3=70\,\mathrm{N},\quad 150^\circ
\]

\[
F_4=50\,\mathrm{N},\quad 270^\circ
\]

Complete la tabla:

\begin{center}
\renewcommand{\arraystretch}{1.8}

\begin{tabular}{|c|c|c|}
\hline
Fuerza & $F_x$ (N) & $F_y$ (N)\\
\hline
$F_1$ & &\\
\hline
$F_2$ & &\\
\hline
$F_3$ & &\\
\hline
$F_4$ & &\\
\hline
$\sum$ & &\\
\hline
\end{tabular}
\end{center}

Determine:

\begin{enumerate}[label=\alph*)]
\item Resultante horizontal.
\item Resultante vertical.
\item Magnitud de la resultante.
\item Dirección y cuadrante.
\end{enumerate}

\espacio{3cm}

\respuesta

\end{ejercicio}

% =========================================
% EJERCICIO CON DIAGRAMA
% =========================================
\newpage

\begin{ejercicio}{15}

Determine la magnitud y dirección
de la fuerza resultante representada
en la figura.

\begin{center}
\begin{tikzpicture}[scale=0.9]

\draw[-{Stealth},gray,thick]
(-4,0)--(4.5,0)
node[right] {$x$};

\draw[-{Stealth},gray,thick]
(0,-3)--(0,3.5)
node[above] {$y$};

\draw[-{Stealth},very thick,blue]
(0,0)--(3,0);

\draw[-{Stealth},very thick,green!60!black]
(0,0)--(2.5,2.1);

\draw[-{Stealth},very thick,purple]
(0,0)--(-2.5,1.8);

\draw[-{Stealth},very thick,red]
(0,0)--(0,-2.5);

\node[below,blue] at (2.2,0)
{$F_1=80\,N$};

\node[above,green!60!black]
at (2.5,2.1) {$F_2=120\,N$};

\node[above,purple]
at (-2.5,1.8) {$F_3=90\,N$};

\node[right,red] at (0,-2.3)
{$F_4=40\,N$};

\draw[orange,thick]
(0.7,0) arc[
start angle=0,
end angle=40,
radius=0.7
];

\node[orange] at (1.2,0.35)
{$40^\circ$};

\draw[orange,thick]
(-0.7,0) arc[
start angle=180,
end angle=145,
radius=0.7
];

\node[orange] at (-1.2,0.3)
{$35^\circ$};

\end{tikzpicture}
\end{center}

\begin{center}
\renewcommand{\arraystretch}{1.7}

\begin{tabular}{|c|c|c|}
\hline
Fuerza & $F_x$ (N) & $F_y$ (N)\\
\hline
$F_1$ & &\\
\hline
$F_2$ & &\\
\hline
$F_3$ & &\\
\hline
$F_4$ & &\\
\hline
$\sum$ & &\\
\hline
\end{tabular}
\end{center}

\espacio{2cm}

\respuesta

\end{ejercicio}

% =========================================
% RETO FINAL
% =========================================
\newpage

\section*{VII. Desafío final}

\begin{reto}

Tres personas tiran de un objeto
mediante cuerdas.

Las fuerzas aplicadas son:

\[
F_1=150\,\mathrm{N}
\quad\text{a }25^\circ
\]

\[
F_2=120\,\mathrm{N}
\quad\text{a }135^\circ
\]

\[
F_3=90\,\mathrm{N}
\quad\text{a }270^\circ
\]

Determine:

\begin{enumerate}
\item Las componentes de cada fuerza.
\item La resultante horizontal.
\item La resultante vertical.
\item La magnitud de la fuerza resultante.
\item Su dirección y cuadrante.
\item La representación gráfica.
\end{enumerate}

\end{reto}

\begin{center}
\plano
\end{center}

\espacio{2cm}

\respuesta

% =========================================
% AUTOEVALUACIÓN
% =========================================
\newpage

\section*{Autoevaluación}

\begin{tcolorbox}[
 colback=celeste,
 colframe=azul,
 title=Revisión de mi aprendizaje
]

Marque las habilidades que
considera que domina.

\vspace{0.3cm}

\casilla\ Identifico magnitudes
escalares y vectoriales.

\vspace{0.4cm}

\casilla\ Represento vectores
gráficamente.

\vspace{0.4cm}

\casilla\ Sumo y resto vectores.

\vspace{0.4cm}

\casilla\ Resuelvo vectores
perpendiculares.

\vspace{0.4cm}

\casilla\ Calculo componentes
rectangulares.

\vspace{0.4cm}

\casilla\ Determino la magnitud
de una resultante.

\vspace{0.4cm}

\casilla\ Determino correctamente
la dirección y el cuadrante.

\end{tcolorbox}

\vspace{0.5cm}

\textbf{¿Qué tema necesito practicar más?}

\vspace{2cm}

\hrulefill

\vspace{0.5cm}

\textbf{Observaciones del docente:}

\vspace{2cm}

\hrulefill

\end{document}
